\MathBookSourcePage{377}
\subsection{习题}
\label{sec:ch13-exercises}
\setcounter{exerciseitem}{0}

\begin{MBExercise}
\label{exr:ch13-cauchy-error-sequence}
证明若对某个 $C<\infty$ 和 $0<\gamma<1$ 有 $\|De^n\|_\Pi\leq C\gamma^n$, 则由~\eqref{eq:ch13-successive-error-relation} 可得 $e^n\to0$. 提示: 在~\eqref{eq:ch13-successive-error-relation} 中取 $v=e^{n+1}-e^n$, 并证明 $\{e^n\}$ 是 $V$ 中的 Cauchy 序列.
\end{MBExercise}

\begin{MBExercise}
\label{exr:ch13-zero-r-iterated-penalty}
证明当 $a(\cdot,\cdot)$ 对称时, 对适当的 $\rho$, 迭代罚方法~\eqref{eq:ch13-iterated-penalty-algorithm} 即使取 $r=0$ 也几何收敛. 提示: 证明对适当的 $\rho$, 存在 $\gamma<1$, 使
\[
\MathBookFormulaID{formula-ch13-ex02-zero-r-contraction}
|a(u,v)-\rho(Du,Dv)_\Pi|^2
\leq\gamma a(u,u)a(v,v).
\]
参见 Glowinski (1984) 的附录 III.
\end{MBExercise}

\begin{MBExercise}
\label{exr:ch13-stokes-iterated-penalty}
假设 $\mathcal T^h$ 满足定理~\ref{thm:ch12-divVh-pressure-errors} 的条件. 令 $V_h$ 如~\eqref{eq:ch12-stokes-velocity-space}, 且 $k\geq4$. 证明当 $r=\rho$ 充分大时, 迭代~\eqref{eq:ch13-iterated-penalty-algorithm} 收敛到 Stokes 方程混合形式~\eqref{eq:ch12-stokes-abstract-problem} 的解. 证明~\eqref{eq:ch13-iterated-penalty-algorithm} 中待解线性系统对称正定. 描述该迭代的停止准则并给出误差估计.
\end{MBExercise}

\begin{MBExercise}
\label{exr:ch13-scalar-elliptic-iterated-penalty}
假设 $\mathcal T^h$ 满足定理~\ref{thm:ch12-divVh-pressure-errors} 的条件. 令 $V_h=V_h^k\times V_h^k$, 且 $k\geq4$. 证明当 $r=\rho$ 充分大时, 迭代~\eqref{eq:ch13-iterated-penalty-algorithm} 收敛到标量椭圆问题混合形式~\eqref{eq:ch12-porous-medium-mixed-problem} 的解. 证明~\eqref{eq:ch13-iterated-penalty-algorithm} 中待解线性系统对称正定. 描述该迭代的停止准则并给出误差估计.
\end{MBExercise}

\begin{MBExercise}
\label{exr:ch13-projected-right-inverse}
证明~\eqref{eq:ch12-operator-L-definition} 蕴含对任意 $p\in\Pi_h$ 有 $P_{\Pi_h}(DLp)=p$. 提示: 回顾 $b(v,q)=(Dv,q)_\Pi$.
\end{MBExercise}

\begin{MBExercise}
\label{exr:ch13-discrete-kernel-projection}
证明 $z\in Z_h$ 等价于 $P_{\Pi_h}(Dz)=0$. 提示: 回顾 $b(v,q)=(Dv,q)_\Pi$.
\end{MBExercise}

\begin{MBExercise}
\label{exr:ch13-augmented-variational-equivalence}
证明~\eqref{eq:ch13-augmented-Lagrangian-matrix-algorithm} 与~\eqref{eq:ch13-augmented-Lagrangian-variational-algorithm} 等价. 提示: 令 $\mathbf W$ 表示 $v$ 关于给定 $V_h$ 基的系数, 并证明
$\mathbf W^t\mathbf B^t\mathbf B\mathbf U^n=b(u^n,P_{\Pi_h}Dv)$.
\end{MBExercise}

\MathBookSourcePage{378}
\begin{MBExercise}
\label{exr:ch13-zero-r-augmented-Lagrangian}
证明当 $a(\cdot,\cdot)$ 对称时, 对适当的 $\rho$, 增广 Lagrange 方法~\eqref{eq:ch13-augmented-Lagrangian-matrix-algorithm} 即使取 $r=0$ 也几何收敛. 提示: 与习题~\ref{exr:ch13-zero-r-iterated-penalty} 类似, 证明对适当的 $\rho$ 存在 $\gamma<1$, 使
\[
\MathBookFormulaID{formula-ch13-ex08-zero-r-augmented-contraction}
|a(u,v)-\rho(P_{\Pi_h}Du,P_{\Pi_h}Dv)_\Pi|^2
\leq\gamma a(u,u)a(v,v).
\]
\end{MBExercise}

\begin{MBExercise}
\label{exr:ch13-Taylor-Hood-augmented-convergence}
令 $V_h$ 如~\eqref{eq:ch12-stokes-velocity-space}, $\Pi_h$ 如~\eqref{eq:ch12-Taylor-Hood-pressure-space}, 且 $k\geq2$. 证明算法~\eqref{eq:ch13-augmented-Lagrangian-variational-algorithm}(或等价地~\eqref{eq:ch13-augmented-Lagrangian-matrix-algorithm})收敛到 Stokes 方程混合形式~\eqref{eq:ch12-stokes-abstract-problem} 的解. 证明~\eqref{eq:ch13-iterated-penalty-algorithm} 中待解线性系统对称正定.
\end{MBExercise}

\begin{MBExercise}
\label{exr:ch13-Taylor-Hood-augmented-explicit}
构造增广 Lagrange 方法, 用于求解时间步进格式~\eqref{eq:ch13-explicit-nonlinear-time-step} 的 Taylor--Hood 离散所对应的方程(参见定理~\ref{thm:ch12-Taylor-Hood-inf-sup-k2}). 陈述并证明相应的收敛结果.
\end{MBExercise}

\begin{MBExercise}
\label{exr:ch13-Taylor-Hood-augmented-semi-implicit}
构造增广 Lagrange 方法, 用于求解时间步进格式~\eqref{eq:ch13-semi-implicit-nonlinear-time-step} 的 Taylor--Hood 离散所对应的方程(参见定理~\ref{thm:ch12-Taylor-Hood-inf-sup-k2}). 陈述并证明相应的收敛结果.
\end{MBExercise}
